% !TeX program = xelatex
\documentclass[11pt,oneside]{article}
\usepackage{ProblemsetStyle}

\hypersetup{
    pdftitle={20260915 Problem set - Solutions},
    pdfauthor={김정우},
    pdfsubject={Science Quiz mathematics practice problem set solutions},
    pdfcreator={XeLaTeX}
}

\begin{document}

\problemsettitle{20260915 Problem set}

\begin{problemsolution}{1}
$x^x>0$이므로 natural logarithm을 취해
\[
    \ln f(x)=x\ln x
\]
를 미분한다. 그러면
\[
    f'(x)=x^x(\ln x+1).
\]
따라서 $f'(x)=0$인 유일한 점은
\[
    x=e^{-1}
\]
이고, 이 점을 기준으로 $f$는 decreasing에서 increasing으로 바뀐다. 그러므로
\[
    \min_{x>0}x^x
    =\left(e^{-1}\right)^{e^{-1}}
    =\boxed{e^{-1/e}}.
\]
핵심은 $x$가 base와 exponent에 동시에 있을 때 logarithmic differentiation을 먼저 떠올리는 것이다.
\end{problemsolution}

\begin{problemsolution}{2}
cyclic group의 homomorphism은 generator의 image로 결정된다.
$\varphi([1]_{12})=[a]_{18}$이라 하면
\[
    12[a]_{18}=[0]_{18}
\]
이어야 하므로
\[
    18\mid12a.
\]
이는 $3\mid a$와 동치이다. 따라서 $\Z/18\Z$에서 가능한 image는
\[
    [0]_{18},[3]_{18},[6]_{18},[9]_{18},[12]_{18},[15]_{18}
\]
의 $6$개이고, 정답은
\[
    \boxed{6}.
\]
일반적으로 additive cyclic groups에 대해
\[
    \bigl|\operatorname{Hom}(\Z/m\Z,\Z/n\Z)\bigr|=\gcd(m,n)
\]
이다.
\end{problemsolution}

\begin{problemsolution}{3}
Stanford Mathematics의 2026년 9월 11일 발표에 따르면, inaugural WLF Prize for Young Scientists에서 K-stability와 birational geometry에 대한 fundamental contributions로 선정된 수학자는 Ziquan Zhuang이다. 따라서 정답은
\[
    \boxed{\text{Ziquan Zhuang}}.
\]
\end{problemsolution}

\begin{problemsolution}{4}
$\mathbf{A}^{\mathsf T}=-\mathbf{A}$이고 matrix size가 odd이므로
\[
    \det(\mathbf{A})
    =\det(\mathbf{A}^{\mathsf T})
    =\det(-\mathbf{A})
    =(-1)^{2027}\det(\mathbf{A})
    =-\det(\mathbf{A}).
\]
따라서 $\det(\mathbf{A})=0$이고
\[
    \rk(\mathbf{A})\leq2026.
\]
한편 $2\times2$ skew-symmetric block
\[
    \begin{pmatrix}
        0&1\\
        -1&0
    \end{pmatrix}
\]
을 $1013$개 두고 마지막에 $0$을 하나 붙인 block diagonal matrix는 rank가 $2026$이다. 따라서 최대 rank는
\[
    \boxed{2026}.
\]
특히 real skew-symmetric matrix의 rank는 항상 even이다.
\end{problemsolution}

\begin{problemsolution}{5}
$f(2026)=q\neq0$이라고 가정하자. $0$과 $q$ 사이에는 irrational number $\alpha$가 존재한다. $f$는 continuous이고
\[
    f(0)=0,
    \qquad
    f(2026)=q
\]
이므로 Intermediate Value Theorem에 의해 어떤 $c$가 존재하여
\[
    f(c)=\alpha
\]
가 되어야 한다. 그러나 $f$의 codomain은 $\Q$이므로 모순이다. 따라서
\[
    \boxed{f(2026)=0}.
\]
사실 같은 argument로 $f$는 $\R$ 전체에서 constant이다. 이는 connected set의 continuous image가 connected라는 사실과도 같다.
\end{problemsolution}

\begin{problemsolution}{6}
Clay Mathematics Institute의 2026 Clay Research Award 발표에서 Burklund, Hahn, Levy, Schlank는 Ravenel의 Telescope Conjecture에 대한 counterexamples를 구성한 업적으로 선정되었다. 따라서 정답은
\[
    \boxed{\text{Ravenel's Telescope Conjecture}}.
\]
이 conjecture는 chromatic homotopy theory의 중심적인 open problem 중 하나였다.
\end{problemsolution}

\newpage

\begin{problemsolution}{7}
$\zeta$가 primitive seventh root of unity이므로
\[
    x^7-1=(x-1)\prod_{k=1}^{6}(x-\zeta^k).
\]
따라서
\[
    1+x+x^2+\cdots+x^6
    =\prod_{k=1}^{6}(x-\zeta^k).
\]
$x=1$을 대입하면
\[
    \prod_{k=1}^{6}(1-\zeta^k)
    =1+1+\cdots+1
    =\boxed{7}.
\]
roots of unity의 product가 보이면 관련 cyclotomic factorization을 먼저 확인하는 것이 핵심이다.
\end{problemsolution}

\begin{problemsolution}{8}
$\{1,2,\ldots,100\}$은 합이 $101$인 $50$개의 pair
\[
    \{1,100\},\{2,99\},\ldots,\{50,51\}
\]
로 partition된다. 합이 $101$인 두 수를 피하려면 각 pair에서 최대 하나만 고를 수 있으므로 최대 $50$개까지 가능하다. 따라서 $51$개를 고르면 Pigeonhole Principle에 의해 어떤 pair에서는 두 수를 모두 고르게 된다. 정답은
\[
    \boxed{51}.
\]
\end{problemsolution}

\begin{problemsolution}{9}
International Mathematical Union의 2026 Fields Medal 수상자 중 Jacob Tsimerman은 Mathematical AI Safety Institute (MAISI)의 Scientific Director를 맡고 있다. 따라서 정답은
\[
    \boxed{\text{Jacob Tsimerman}}.
\]
MAISI는 powerful AI systems의 safety를 위한 mathematical foundations를 연구하는 independent nonprofit이다.
\end{problemsolution}

\begin{problemsolution}{10}
두 수를
\[
    A=2^{2026}-1,
    \qquad
    B=2^{2024}-1
\]
이라 두자. Euclidean algorithm의 한 단계로
\[
    A-4B
    =(2^{2026}-1)-4(2^{2024}-1)
    =3.
\]
따라서 $\gcd(A,B)$는 $3$의 divisor이다. 한편 $2026$과 $2024$는 모두 even이므로
\[
    2^{2026}\equiv2^{2024}\equiv1\pmod3,
\]
즉 $3\mid A$이고 $3\mid B$이다. 따라서
\[
    \boxed{\gcd\!\left(2^{2026}-1,2^{2024}-1\right)=3}.
\]
일반적으로
\[
    \gcd(a^m-1,a^n-1)=a^{\gcd(m,n)}-1
\]
이라는 identity가 성립하며, 여기서는 $\gcd(2026,2024)=2$이다.
\end{problemsolution}

\end{document}
